% sec-04: what an external standard requires.  Table 9, Figure 6 (d2-type grid).
\section{What an External Standard Requires}

A request to be listed as an external standard invites one question before any
other: a standard by what measure? This section proposes eight criteria that any
body of work would have to meet before a reviewer could reference it as a
benchmark for oncology clinical trial research, and scores the company's record
against each. The criteria draw on the ASME approach to computational model
credibility \cite{asmevv40}, the ICH principles for model-informed drug
development \cite{ichm15}, and the FDA's draft guidance on artificial
intelligence in regulatory decision-making \cite{fdaaiguidance2025}.

Three of the eight are met across the record: persistent identifiers, dated
versions, and an open method. The other five are met in part or not yet, and
each is taken in turn below before Table 9 summarizes all eight.

\begin{appfloat}[!tb]
\begin{appfig}[d2-type / grid with shaded status cells]
% Rows: eight criteria.  Columns: five parts of the record.  Cell fill encodes
% status: accent = carries it, mid = in part, light gray = does not.
\foreach \x/\lab in {3.30/{2025 papers}, 5.30/{Protocols and IND}, 7.30/{Applications and plans}, 9.30/{Repository and builds}, 11.30/{Outside replies}}{
  \node[d2cellh,minimum width=19mm,text width=17mm,minimum height=8mm] at (\x,0) {\lab};}
\node[d2cellh,minimum width=38mm,text width=36mm,minimum height=8mm] at (0.60,0) {Criterion};
\foreach \y/\c in {-0.78/{1. Persistent identifiers}, -1.46/{2. Dated versions}, -2.14/{3. Open method}, -2.82/{4. Credibility assessment}, -3.50/{5. Regulatory mapping}, -4.18/{6. Independent review}, -4.86/{7. Outcome validation}, -5.54/{8. Adoption, confirmed}}{
  \node[d2celll,minimum width=38mm,text width=36mm,minimum height=6.8mm] at (0.60,\y) {\c};}
% one row per criterion: fills for the five columns, left to right
\foreach \y/\a/\b/\c/\d/\e in {%
  -0.78/fundaccent/fundaccent/fundaccent/fundmid/fundgrayl,
  -1.46/fundaccent/fundaccent/fundaccent/fundaccent/fundgrayl,
  -2.14/fundaccent/fundmid/fundmid/fundaccent/fundgrayl,
  -2.82/fundaccent/fundgrayl/fundgrayl/fundmid/fundgrayl,
  -3.50/fundgrayl/fundaccent/fundmid/fundmid/fundgrayl,
  -4.18/fundmid/fundmid/fundmid/fundmid/fundgrayl,
  -4.86/fundmid/fundgrayl/fundgrayl/fundgrayl/fundgrayl,
  -5.54/fundgrayl/fundgrayl/fundgrayl/fundgrayl/fundgrayl}{
  \foreach \x/\f in {3.30/\a, 5.30/\b, 7.30/\c, 9.30/\d, 11.30/\e}{
    \node[d2cell,fill=\f,minimum width=19mm,minimum height=6.8mm] at (\x,\y) {};}}
\legkey{-1.30}{-6.35}{fundaccent}{Carries the criterion}
\legkey{2.90}{-6.35}{fundmid}{Carries it in part}
\legkey{6.60}{-6.35}{fundgrayl}{Does not carry it yet}
\ptitle{-1.30}{0.75}{Where each criterion is carried}
\end{appfig}
\vspace{-0.60cm}
\figcaption{Figure 6. The eight criteria against five parts of the record, shaded by whether each\\
part carries the criterion fully, in part, or not yet; the last two rows are open.}
\end{appfloat}

\subsection{Where each criterion is carried}

Before the criteria are taken one by one, Figure 6 shows where in the record each
is carried. It breaks the record into five parts and shades each cell by whether
that part carries the criterion fully, in part, or not at all. Two things stand
out. The first three criteria are carried across the whole record. The last two
are carried only by replies the company does not yet have.

\subsection{The five criteria not yet fully met}

\textbf{Credibility assessment.} The score of 81.9 over 55 tests is a pre-trial
credibility score for the digital twin, not a post-trial validation, and it
covers the model rather than the trial documents \cite{daraxsims}. Extending the
same test battery to the Phase 1 protocol's decision rules is the next step.

\textbf{Regulatory mapping.} The investigational new drug application and both
protocols are mapped to 21 CFR Part 312 and adapted ICH guidance
\cite{phase1ind,ecfr2026part312}. No regulator has reviewed the mapping, and the
routing letter to the FDA program asks where such a review would belong.

\textbf{Independent review.} Each document is produced by one model and reviewed
by two others from different developers \cite{tripartisan2026}. A clinical
investigator's written review of the Phase 1 protocol, published with consent,
would add the human reviewer a standard needs.

\textbf{Outcome validation.} The only comparison between a company prediction
and a trial result is a chronology. A prospective test means publishing a dated
prediction for one endpoint of a trial that has not yet reported, and being
scored on it afterward.

\textbf{Adoption.} The company can document its own use of OpenAI and Anthropic
tools. It cannot document either company's internal use of its papers, and it
has written to both inviting a correction. Until either replies, the criterion
is unmet on independent evidence, and the company says so.

\subsection{The eight criteria together}

\begin{apptable}
\begin{tabularx}{\textwidth}{>{\raggedright\arraybackslash}p{0.5cm} >{\raggedright\arraybackslash}p{4.4cm} >{\raggedright\arraybackslash}X >{\raggedright\arraybackslash}p{1.6cm}}
\headrow \thead{No} & \thead{Criterion} & \thead{What the record holds today} & \thead{Status}\\
\midrule
1 & Persistent identifiers & Eleven DOIs across the program & Met \\
2 & Dated versions & Deposit dates; releases v4.0.0 to v4.9.0 & Met \\
3 & Open method & Public source, prompts and build records & Met \\
4 & Credibility assessment & VVUQ score 81.9 over 55 tests, pre-trial & In part \\
5 & Regulatory mapping & 21 CFR 312 and ICH mapping, unreviewed & In part \\
6 & Independent review & Two model reviewers; no human reviewer & In part \\
7 & Outcome validation & One chronology; no prospective test & Not yet \\
8 & Adoption, confirmed & The company's statement, unconfirmed & Not yet \\
\end{tabularx}
\end{apptable}
\vspace{-0.60cm}
\tabcap{Table 9. Eight criteria for an external standard scored against the record as\\
it stands today: three are met, three are met in part, and two are not yet met.}

A table that scored every criterion as met would be read as advocacy. This one
says two are not yet met, and the subsection above names the action that would
change each of the five that fall short. That is what makes it a record rather
than a pitch.
